% Figure: Gibbs phenomenon. Fourier sine partial sums of the
% indicator function on (0, 1/2) overshoot the jump by ~9% of the
% jump height, with the overshoot concentrating but not vanishing
% as more terms are added.
\begin{figure}[htbp]
\centering
\begin{tikzpicture}
\begin{axis}[
  width=12cm, height=6.5cm,
  xlabel={position $x$},
  ylabel={partial sum $S_N(x)$},
  xmin=0, xmax=1,
  ymin=-0.25, ymax=1.25,
  grid=both,
  major grid style={engblack!20},
  minor grid style={engblack!10},
  legend pos=north east,
  legend cell align=left,
  font=\small,
  samples=400,
  domain=0.0005:0.9995,
]
% Target: square pulse = 1 on (0, 0.5), 0 on (0.5, 1).
\addplot[engblack, very thick, dashed, samples=4] coordinates
  {(0,1) (0.5,1) (0.5,0) (1,0)};
\addlegendentry{target (jump at $x=1/2$)}

% Fourier sine coefficients of indicator of (0,1/2) on (0,1):
% b_n = (2/L) int_0^{1/2} sin(n pi x) dx = (2/(n pi))(1 - cos(n pi/2)).
% Partial sum N=5.
\addplot[engblue, thick] {
    (2/pi)*( (1-cos(deg(pi/2)))*sin(deg(pi*x))/1
    + (1-cos(deg(2*pi/2)))*sin(deg(2*pi*x))/2
    + (1-cos(deg(3*pi/2)))*sin(deg(3*pi*x))/3
    + (1-cos(deg(4*pi/2)))*sin(deg(4*pi*x))/4
    + (1-cos(deg(5*pi/2)))*sin(deg(5*pi*x))/5 )
};
\addlegendentry{$N = 5$}

% Partial sum N=21 (sum of nonzero odd-ish terms; include all up to 21).
\addplot[engred, thick] {
    (2/pi)*(
      (1-cos(deg(pi/2)))*sin(deg(pi*x))/1
    + (1-cos(deg(2*pi/2)))*sin(deg(2*pi*x))/2
    + (1-cos(deg(3*pi/2)))*sin(deg(3*pi*x))/3
    + (1-cos(deg(4*pi/2)))*sin(deg(4*pi*x))/4
    + (1-cos(deg(5*pi/2)))*sin(deg(5*pi*x))/5
    + (1-cos(deg(6*pi/2)))*sin(deg(6*pi*x))/6
    + (1-cos(deg(7*pi/2)))*sin(deg(7*pi*x))/7
    + (1-cos(deg(8*pi/2)))*sin(deg(8*pi*x))/8
    + (1-cos(deg(9*pi/2)))*sin(deg(9*pi*x))/9
    + (1-cos(deg(10*pi/2)))*sin(deg(10*pi*x))/10
    + (1-cos(deg(11*pi/2)))*sin(deg(11*pi*x))/11
    + (1-cos(deg(12*pi/2)))*sin(deg(12*pi*x))/12
    + (1-cos(deg(13*pi/2)))*sin(deg(13*pi*x))/13
    + (1-cos(deg(14*pi/2)))*sin(deg(14*pi*x))/14
    + (1-cos(deg(15*pi/2)))*sin(deg(15*pi*x))/15
    + (1-cos(deg(16*pi/2)))*sin(deg(16*pi*x))/16
    + (1-cos(deg(17*pi/2)))*sin(deg(17*pi*x))/17
    + (1-cos(deg(18*pi/2)))*sin(deg(18*pi*x))/18
    + (1-cos(deg(19*pi/2)))*sin(deg(19*pi*x))/19
    + (1-cos(deg(20*pi/2)))*sin(deg(20*pi*x))/20
    + (1-cos(deg(21*pi/2)))*sin(deg(21*pi*x))/21 )
};
\addlegendentry{$N = 21$}

% Reference overshoot level ~1.09 (Gibbs constant).
\addplot[engblack, thin, dotted, samples=2, domain=0:1] {1.0895};
\node[engblack, font=\scriptsize, anchor=south east] at (axis cs:0.49,1.09)
  {$\approx 9\%$ overshoot};
\end{axis}
\end{tikzpicture}
\caption[Gibbs phenomenon at a jump discontinuity]{Fourier sine
partial sums $S_N(x)$ for the indicator function that equals one on
$(0,1/2)$ and zero on $(1/2,1)$. As $N$ grows from 5 to 21 the
approximation improves in the bulk but the overshoot at the jump
does not shrink: it settles at roughly $9\%$ of the jump height (the
Gibbs constant, dotted line) and concentrates into an ever-narrower
band against the discontinuity at $x = 1/2$. The series converges to
the midpoint value $1/2$ at the jump itself. The engineer who fits a
truncated Fourier expansion to discontinuous data should expect this
ringing near sharp features and decide in advance whether the
application tolerates it or calls for a different basis.}
\label{fig:vol02:ch12:gibbs}
\end{figure}
